\section{Lacunas de Pesquisa: A Necessidade de Novos Paradigmas Metodológicos}

O diagnóstico das limitações anteriormente expostas converge para um sintoma central: a agenda de pesquisa atual em odontologia digital encontra-se saturada de abordagens reducionistas, sendo imperativa a transição para métodos de pesquisa mais robustos e sistêmicos.

\subsection{O Desafio dos Estudos Longitudinais}

A quase totalidade do corpus científico sobre gêmeos digitais odontológicos analisa janelas temporais agudas — tipicamente limitadas ao período do tratamento ativo ou a seguimentos não superiores a 24 meses. Contudo, em disciplinas reconstrutivas como a implantodontia e a reabilitação oral, complicações biológicas e mecânicas sérias (como peri-implantite, afrouxamento de parafusos, fraturas de zircônia ou perda de crista óssea) frequentemente se manifestam na marca de 5 a 10 anos. Não há, até o momento, ECRs de longo prazo que comparem os índices de sobrevida de reabilitações planejadas exclusivamente com gêmeos digitais preditivos versus métodos de rotina \cite{elsaid2025digital}. Sem a validação longitudinal, a jactada ``previsibilidade absoluta'' atribuída à tecnologia é uma construção teórica, e não um fato científico.

\subsection{Viés Étnico e Demográfico (\textit{Data Bias})}

Existe um profundo viés geográfico e fenotípico que permeia a literatura global: os coeficientes biomecânicos, os tensores de elasticidade, as espessuras corticais e as morfometrias mandibulares que alimentam as bibliotecas dos softwares de gêmeos digitais são, majoritariamente, derivados de amostras caucasianas norte-americanas ou asiáticas \cite{khoshfekr2025digital}. A exportação acrítica desses parâmetros (\textit{hardcoding}) para a complexa matriz de miscigenação da população brasileira pode gerar previsões de movimentação ortodôntica ou distribuição de tensões perigosamente equivocadas. A pesquisa nacional precisa focar no mapeamento ontológico da biometria craniofacial brasileira, calibrando modelos com dados locais.
